\section{Modified Principal Nest}\label{principal nest}
                      

\subsection{Generalized polynomial-like maps}\label{GPL}


 A {\it generalized polynomial-like map} (GPL map) is a holomorphic map
$g\from \cup W_i \to V$, where $V \subset \C$ is a topological disk and 
$W_i\Subset V$ are topological disks with disjoint closures such that the restrictions
$g:W_i\ra V$ are branched coverings, and moreover, all but finitely many of them have degree one.

\ssk {\it Remark.} To prove Theorem B in full generality, we need to allow infinitely many disks
$W_i$. However, in the ``persistently recurrent'' case that interest us most
it is enough to consider GPL maps defined on finitely many disks $W_i$.
\ssk

\comm{********
 A {\it generalized polynomial-like map} (GPL map) is a proper holomorphic map
$g\from \W \to V$, where $V \subset \C$ is a topological disk and $\W \Subset V$ is full.
Then $\W = \bigcup_{i=0}^k W_i$, where the domains $W_i$ are topological disks
 and the map $g: W_i \ra V$ is a branched covering. 
***************}

We let $K_g  = \bigcap_{n=0}^\infty g^{-n} V$ be the set of points of $V$ on which $g$ is infinitely iterable
(the ``filled Julia set'').


% By shrinking $V$ slightly we can assume that it is a quasidisk with piecewise smooth boundary, 
% and that the components  $W_i$ have disjoint closures. 
% We will always assume this in what follows.
%
%\ssk {\it Remark.} In fact $\di V$ can be assumed to be real analytic but it is often convenient to 
%    allow a few corners on it. 
%\ssk 

A GPL map $g$ is called {\it unicritical} if it has a single critical point.
{\it In what follows we will consider only unicritical  GPL maps, and we will always put its critical point at $0$.} 
Let $d$ be  the local degree of $g$ near 0. We let   $W_0\equiv W$ be the ``central domain'', that is, the one
containing~0.  

The {\it postcritical set} $\OO_g$ of a (unicritical) GPL map is the closure of the orbit $\{g^n 0\}_{n=0}^\infty$.
% The domains $W_i$ that do not intersect $\OO_g$ are not essential for dynamics,
% so in what follows we consider only the maps for which all the domains $W_i$ do intersect $\OO_g$. 


\msk
  {\it Puzzle pieces} of depth $n$ of a GPL map $g$ are components of $g^{-n}(V)$. Puzzle pieces containing
$0$ are called {\it critical}. 

If the critical point returns to some critical puzzle piece $A$,
then  the first return map $h$ to $A$ is also GPL. Let   $\cup B_i$ be its domain of definition.  
Restricting $h$ to the union of those components $B_i$ that intersect the postcritical set,
we obtain a GPL map called the {\it generalized renormalization} 
$r_A(g)$ of $g$ on $A$. 

If we do not specify the domain $A$ of the generalized renormalization,
then it is assumed to be $W$, so $r(g)\equiv r_W(g)$. 


% \begin{lem}
%  Let $f: z\mapsto z^d+ c$ be a non-renormalizable map with all fixed points repelling.
%If the postcritical set $\OO_f$ is not minimal then the Julia set $J(f)$ is locally connected. 
%\end{lem}
                                                                                                   

\subsection{Dynasty of kingdoms}\label{dynasty}
Let us introduce a modified notion of (unicritical) GPL map called a kingdom map. 

Let us consider three topological disks,  $W\supset U\Supset  A\ni 0$,
called the {\it kingdom domain}, the {\it castle},  and the {\it king} respectively. 
Let $D_j\Subset W\sm\bar A$ be a family of topological disks (`` king's subjects'')
such that $\bar D_j\cap \di U=\emptyset$.
Finally, let  $M_k\Subset U\sm\bar A$ be another family of topological disks 
(``king's men'').
A map
$$
     G: A \cup D_j\cup M_k \ra W
$$
is called a {\it kingdom map} (of {\it local degree} $d$) if
\begin{itemize}
 \item  The closures $\bar A$, $\bar D_j$ and $\bar M_k$ are pairwise disjoint; 
 \item $G: A\ra W$ is  a $d$-to-$1$ branched covering ramified only at 0;
 \item Each $G: D_j\ra W$  is a biholomorphic isomorphism; 
 \item Each $G: M_k\ra U$ is is a biholomorphic isomorphism. 
\end{itemize}
We let $\OO_G$ be the postcritical set of the kingdom map $G$. 


\begin{figure}[htbp]
\begin{center}
\input{kingdom.pstex_t}
\caption{Kingdom map}
\end{center}
\end{figure}


% \makefigure{kingdom}{Kingdom map}

When $U=W$, kingdom maps become GPL maps.

\msk
Let us now consider a (unicritical)  GPL map $g: \cup W_i\ra V$, $W\equiv W_0$. 
Let us define  the {\it kingdom renormalization} $G=R(g)$ of $g$
whose result will be a kingdom map $G$. 

 
 If $g(0)\in W$ then we say that the {\it central return} occurs.
If $g^k (0)\in W$ for $k=0, \dots, N-1$ but $g^N(0)\not\in W$, 
then we have a nest of topological disks
\begin{equation}\label{cascade}
    V\equiv \Om_0\Supset W\equiv  \Om_1\Supset\dots \Supset \Om_N\equiv U
\end{equation}
such that $g: \Om_{k+1}\ra \Om_k$ is a unicritical branched covering of degree $d$ 
% (with the only branched point at $0$)
and $g(0)\in \Om_{N-1}\sm \Om_N$.   
This nest is called a {\it central cascade (of length $N$)}. 
Note that the non-central-return event corresponds to the cascade of length 1. 

In  the kingdom renormalization $Rg$,  $W$ will be the kingdom domain and $U$ will be the  castle.
%and the annuli $A^k = \Om_{k-1} \sm \Om_k$ are {\it defense belts}.   

Let us consider the first return map to $W=\Om_1$:
$$
   h : X_0 \cup \bigcup_{i>0} X_i  \ra W, 
$$
where $X_0\ni 0$ (in case when $N\geq 2$, $X_0= \Om_2$). 
% $h_k|\Om^k = g|\Om^k$  and the $\Om^k_i$, $i>0$,   are topological disks compactly contained in $A^k$. 
% Those domains $\Om^k_i$ that intersect the postcritical set $\OO_g$ are called {\it king's subjects}.   

Let us consider the domain $X_s$, $s>0$,  containing  $g^{N-1}(0)$. 
Then  the pullback $A  = g^{-(N-1)}(X_s)|\, 0$ is the  king of $Rg$.  
The kingdom map $G$ on $A$ is defined as $ h \circ g^{N-1}: A \ra W$.
Notice that it is a unicritical $d$-to-$1$ branched covering. 

Let us define {\it king's subjects} $D_j$ as non-critical pullbacks of the domains $X_i$ ($i\not= 0$) under the maps 
$$
   g^{k-1}: \Om_k\ra W,\quad k=1,2,\dots, N,
$$
that intersect the postcritical set. 
Thus,  each subject $D_j$ is  univalently mapped  onto some $X_i$, $i>0$,
by an appropriate map $g^{k-1}: \Om_k \ra W$, $k\in [1,N]$.
On this subject let us define the kingdom  map $G: D_j\ra W$ as $G|D_j = h\circ g^{k-1}|D_j $.
Obviously, it is a biholomorphic isomorphism. 

% If the postcritical set  $\OO_g$ is minimal then the king has only finitely many subjects. 

Finally, let us define  king's men $M_k$ 
as the pullbacks of $U$ under  $g: U\ra g(U)\supset U$ that intersect the postcritical set.
There are at most $d$ king's men, and $g$ univalently maps each of them onto $U$. 
Let $G|M_k= g|M_k$. 

Thus, we have defined the desired kingdom renormalization 
$$
   G=R(g) : A \cup D_j\cup M_k \ra W.
$$ 

Let $g^N(0)\in W_j$, $j>0$. For $G=R(g)$,  
let us define the {\it king's apartment} $\De$ as $g^{-N}(W_j)|\, 0$. 
Then $A\Subset \De\subset U$
and the map $g^{N+1} : \De\ra V$ is a unicritical $d$-to-1 branched covering. 
% (ramified only at $0$).
This creates a collar $\De\sm A$ around the king.  

\ssk
{\it Remark.} If $N=1$ (i.e., the non-central return occurs under $g$), then the kingdom renormalization
$G=R(g)$ coincides with the generalized renormalization defined in \S \ref{GPL}. 

%The above data (that is, the central cascade, the king in his apartment,  with his subjects and men,
%ruled by the kingdom map $G$)  will be called the {\it kingdom  associated with $g$}.     
%Somewhat abusing notations, we will denote the whole kingdom by $G$.  

\msk
 Given a kingdom map $G$, let us define its renormalization $g=r(G)$ as the first return map
$g: \cup B_i\ra A$  to the king $A$ restricted to those domains $B_i$ that  intersect the postcritical set $\OO_G$. 
It is a unicritical GPL map.
%% If  $\OO_G$ is minimal then there are only finitely many domains $B_i$,  so that  $g$ a GPL map. 

%Notice that on the postcritical set $\OO_g$, 
%the first return map $h$ to $K$ by $g$ and  $G$ are the same.
%It extends to a quadratic-like map $h: \cup \De_j\ra K$.            \note{if $\OO_g$ is minimal} 
%The kingdom map associated to $h$ is called the {\it kingdom renormalization}, $R(G)$, of $G$.

Beginning with some GPL map $g\equiv g_0$, 
%  with minimal postcritical set,
we construct in the above way  a {\it dynasty} of kingdoms, that is, 
 a sequence  $(g_n, G_n)$ such that $g_n$ is a GPL map, $G_n$ is a kingdom map,
$G_n= R(g_n)$ and $g_{n+1} = r(G_n)$. This dynasty terminates if and only if:

\ssk\nin $\bullet$ The map $g$ is combinatorially non-recurrent, that is, 
   the critical point does not return to some critical puzzle piece; or 

\ssk\nin $\bullet$
 Some map $g_n$ has an infinite central cascade, i.e.,
 it is a Douady-Hubbard polynomial-like map \cite{DH}
with non-escaping critical point.
In this case $g$ is called {\it renormalizable in the sense of Douady and Hubbard}.  

\ssk
When we consider a dynasty of kingdoms $(g_n, G_n)$,  the associated domains will be marked with
superscript $n$ (e.g., $V^n$, $W^n$, etc.) 
However, we usually skip the label when we are concerned with a single kingdom.

\ssk {\it Remark.} 
  It is easy to see that the maps $g_n$ coincide with the generalized renormalizations of $g$ on
domains $V^n$ as defined in \S \ref{GPL}, i.e., $g_n= r_{V^n} (g)$. 


%The level $n+1$ is called {\it buffer} if $U^n\not= V^n$ and the return $G_n: V^{n+1} \ra V^n$ is central, i.e.,
%$G_n(0)\in V^{n+1}$. 
%Note that there are no two consecutive buffer levels. 

The nest 
$$
  V^0\supset W^0\supset \dots\supset    W^{n-1} \supset   V^n \supset W^n \supset V^{n+1}\supset\dots
$$
is called a {\it Modified Principal Nest}. 
Sometimes it is convenient to relabel it  in a uniform way: 
\begin{equation}\label{E}
  E^0\supset E^1\supset  \dots \supset  E^i \supset E^{i+1} \supset\dots,
\end{equation}
so that $V^n=E^{2n}$, $W^n= E^{2n+1}$. 
The consecutive $E$-domains are dynamically related: $E^{i-1}= \psi_i(E^i)$, 
where $\psi_i$ is a unicritical $d$-to-1 branched covering which is an appropriate iterate of $g$.

  

\subsection{First king}\label{sec: first king}
We will describe in this section how to  associate to a unicritical polynomial $f:z\mapsto z^d+c $ 
(or, more generally, polynomial-like map) a  dynasty of kingdom maps.   
Our standing assumption is that {\it the Julia set $J_f$ is connected and all periodic points of $f$ are
repelling}.  Then $f$ has $d-1$ {\it non-dividing} fixed points $\beta_i$ (landing points of the external rays%
\footnote{In the case of polynomial-like map, external rays are defined by means of straightening.}
with angles $2\pi/(d-1)$), and one {\it dividing} fixed point $\alpha$. There are $q>1$ external rays $\RR_i^0$ landing at 
$\alpha$ which are cyclically permuted by the dynamics, see \cite{M-rays}.

 
Let us select some equipotential $E^0$; it bounds some topological disk $Q^0$.  
The rays $\RR_i^0$ divide $Q^0$ into $q$ disks $Y^0_i$ called the {\it Yoccoz puzzle pieces} of depth 0.  
Let $Y^0\equiv Y^0_0$ stand for the critical puzzle piece, i.e., the one containing 0. 

The equipotential  $E^1=f^{-1}E^0$ bounds some topological disk $Q^1$.
Let us consider  $dq$  rays of  $f^{-1} (\cup\RR_i^0)$. 
They divide $Q^1$ into $(q-1)d+1$ topological disks called Yoccoz puzzle pieces of depth 1. 
Let $Y^1$ stand for the critical puzzle piece of depth 1. 
There are also $q-1$  puzzle pieces $Y^1_i$ of depth 1 contained in the corresponding off-critical pieces of depth 0.
All other puzzle pieces of depth 1 will be denoted $Z_j^1$. They are attached to the $f$-preimages of $\alpha$
that are different from $\alpha$ itself. 
% \note{picture} 

  The map $f$ is called   {\it satellite renormalizable}
(or, {\it immediately renormalizable})  if
 $$
           f^{lq}(0)\in Y^0,\quad  l=0,1,2\dots.
$$
In this case,  we let $Y^{lq}  = f^{-lq}(Y^0)|\,0$ and consider the unicritical branched covering
$f^q: Y^q\ra Y^0$ of degree $d$. By slight ``thickening'' of the domain of this map (see \cite{M}), 
it can be turned into a unicritical GPL map called the {\it (satellite) renormalization} $Rf$ of $f$.%
\footnote{In the context of GPL maps we use the term ``DH renormalization''
to distinguish it from the generalized renormalization.
In the polynomial case, we refer to it as just ``renormalization'',
as it should not lead to confusion.} 

\ssk
In the satellite renormalizable case, $f$ does not originate any dynasty. 
Otherwise, 
there exists an $l\in \N$ such that $f^{lq}(0)$ belongs
to some puzzle piece $Z_j^1$.
In this case, we let $V^0= f^{-lq}(Z_j^1)|0$ be the {\it first kingdom}, 
and we let $g\equiv g_0: \cup W_i^0\ra V^0$ be the first return map to $V^0$. 
It is easy to check that  $W_i^0 \Subset V^0$. 
%
% If the postcritical set $\OO_f$ is minimal then only finitely many of the domains $W_i$ intersect it.
% Restricting $\tl g_0$ to the union of those domains $W_i$ that do not intersect $\OO_f$, 
% we obtain the first map $g\equiv g_0$ in the dynasty. 
Let $G_0$ be the associated kingdom map. It originates the dynasty $(g_n, G_n) $ associated with $f$.


\ssk
The map $f$ is called {\it primitively renormalizable} if its dynasty contains  a quadratic-like map $g_n: W^n\ra V^n$
with connected Julia set. This quadratic-like map  is called the {\it (primitive) renormalization} $Rf$ of $f$.
In this case, we cannot construct the next kingdom map  $G_n$,
so the dynasty terminates. 
It also terminates if the map $g$ is combinatorially non-recurrent. 
Otherwise, the process can be continued indefinitely,
and the dynasty $(g_n, G_n)$ is eternal.

If the map $f$ is renormalizable (either in the satellite or in the  primitive sense), 
we can take its renormalization $Rf$
and consider its dynasty. If $Rf$ is renormalizable, 
we can pass to the second renormalization $R^2 f$,   and so on. 
If the map $f$ is at most finitely renormalizable, 
in the end we obtain a non-renormalizable quadratic-like map $R^m f$.
% It is either combinatorially non-recurrent or it has the eternal dynasty. 
This is the map we will be working with. 
So, {\it in what follows we will assume that $f$ itself is non-renormalizable.}

From now on, we can forget about the original polynomial $f: z\mapsto z^d +c $,
and  replace it with the first map $g: \cup W_i^0\ra V^0$ of the associated dynasty.


\subsection{Extensions}

Let us begin with a trivial but  useful observation:

\begin{lem}[Telescope]\label{telescope}
  Let $X_k$ be a sequence of topological disks, $k=0,1,\dots, m$,
and let $\phi_k : X_k \ra \phi(X^k)$  be  branched coverings of degree $d_k$ such that
$\phi(X_k)\supset X_{k+1}$. Let $\Phi= \phi_{n-1}\circ\dots\circ \phi_0$ (wherever it is defined),
and let $P\subset X_0$ be a component of its domain of definition.
Then $\Phi: P \ra V_n$ is a branched covering of degree at most $d_0\cdot\cdot\cdot d_{n-1}$.
\end{lem}


\begin{lem}\label{landing extension}
  Let $g^m z\in A $ be the first landing  of the $\orb(z)$ at $A$.
Then there exists a puzzle piece $P\ni z$ such that $g^m$ univalently maps $P$ onto $U $. 
% Moreover, the puzzle pieces $g^k P$, $k=0,\dots, m-1$, are pairwise disjoint. 
\end{lem}  

\begin{pf}
  Let $P= g^{-m} (U) |z$. Then $g^m = h^k \circ g^s$, where $g^s(z)$ is the first landing of $\orb z$
at $U$, $h=h_U: \cup B_i\ra U$ is the generalized renormalization on $U$, and $k$ is the first landing moment
of $\orb_h(g^s(z))$ at $A$. It is a simple exercise to show that $g^s$ is univalent
on $g^{-s}(U)|z$. Moreover,  $h$  univalently maps each non-central component $B_i$, $i>0$, onto $U$.
Now the first assertion follows from the Telescope Lemma.


\end{pf}  


\begin{cor}\label{return extension}
  Let $z\in A$, and let $g^m z\in A$ be the first return of the $\orb(z)$ to $A$.
Let  $P = g^{-m}(A)| z$.                     
% Then $P \subset A$. 
If  $P$ is not critical then the  map  $g^m: P \ra A$   is univalent.
 Otherwise  $g^m: P \ra A$ is  a unicritical branched covering of  degree $d$. 
\end{cor}

\begin{pf}
 Decompose  $g^m: P\ra A$ as $g: P\ra g(P)$ and the first landing map $g^{m-1}: g(P)\ra A$.
\end{pf}

\comm{***********************************************************
\ssk\nin {\rm (ii)}
  Let  $z=0$, and let $D$ be any component of $g^{-m}(U)$. % contained in $A$. 
Then  $g^m: D\ra U$ has at most one critical point. 
\end{cor}  

\begin{pf}
(i) The first assertion follows from the fact that at the first moment $T$ when $g^T A=W\ni 0$,
    we have: $g^T A \supset U$. 
    The rest follows from  Lemma \ref{landing extension} by decomposing $g^m|\sm P$ as $L\circ g$,
    where $L$ is the first landing map on $g(P)$. 

\ssk
(ii)
Consider the orbit of $D$,   
$$
        D_k= g^k D, \quad k=0,1,\dots, m-1.
$$
If there are no critical pieces in it, then the map $g^m: D\ra U$ is univalent.
Otherwise, let $D_s$ be the last critical piece in this orbit.  Then $g^{m-s}: D_s\ra U$ is unicritical. 
Moreover, by (i) $D_s\subset A$, and hence the first return time of the critical point  to $D_s$ is at least $m$.
It follows that the preceding  puzzle pieces, $D_0,\dots, D_{s-1}$, are off-critical,
 and the conclusion follows.   
\end{pf}
**********************************************************}


Applying this to the first return of the critical point to $A^{n-1}=V^n$, we obtain:

\begin{cor}\label{space around W}
   The map $g_n: W^n \ra V^n$ admits an analytic extension to a puzzle piece $\tl W^n\Supset W^n$
such that $g_n: \tl W^n \ra U^{n-1}$ is a unicritical branched covering of degree $d$.  
Moreover, $\tl W^n\subset V^n$.
% Moreover, letting $g_n|\, W^n =g^{t_n}|\, W^n $, the map $g^{t_n}$ has at most one critical point on 
% every component of $g^{-t_n}(U^{n-1})$.
\end{cor}



Let us now construct similar extensions for kingdom maps: 

\begin{lem}\label{cascade extension}
 There is puzzle piece $\tl A \Supset A$ such that  
the map $G: A \ra W $  admits a unicritical degree $d$  extension to a map $\tl A  \ra V$.
Moreover, $\tl A\subset \De$ where $\De$ is the king's apartment.
\end{lem}

\begin{pf}
The map $G: A\ra W$ can be decomposed as $g^k\circ  g^N$ where $N$ is the length of the central cascade of $g$,
% $g:\cup W_i\ra V$ is the generalized renormalization on $V$,
 and $k$ is the first entry time of $\orb_g (g^N(A))$ to $W$ (recall that $g^N(A)\subset V\sm W$). 
The map $g^N: A\ra g^N(A)$ admits an analytic  extension to  a unicritical $d$-to-1 branched covering
$g^N: \De\ra W_j$ for some $j>0$, while $g: W_i\ra V$ is a biholomorphic isomorphism for any $i>0$.
Now the conclusion follows by the Telescope Lemma. 
\end{pf}

% Sometimes it will be convenient to relable the nest 
%$$
%  \dots\supset    W^{n-1} \supset   V^n \supset W^n \supset V^{n+1}\supset\dots
%$$
%in a uniform way: 
%\begin{equation}\label{E}
%    \dots \supset  E^i \supset E^{i+1} \supset\dots,
%\end{equation}
%so that $V^n=E^{2n}$, $W^n= E^{2n+1}$. 
%The consecutive $E$-domains are dynamically related: $E^{i-1}= \psi_i(E^i)$, 
%where $\psi_i$ is a $d$-to-1 branched covering which is an appropriate iterate of $g$.

Let us define {\it enlargements} $\hat E^i$ of domains $E^i$
of the Modified Principal Nest (\ref{E})  as follows: 
$\hat W^n =V^n$ and $\hat V^n = \De^{n-1}$.
We also have the {\it buffers} $\tl E^i\subset \hat E^i$
 constructed in Corollary  \ref{space around W} and Lemma \ref{cascade extension}.
These lemmas tell us that  any map $\psi_i$ analytically extends to a unicritical  $d$-to-1 branched covering 
$\psi_i: \tl E^i\ra \hat E^{i-1}$. 
 For $i<k$, let 
$$
  \Phi_{i,k} = \psi_{i+1}\circ\dots\circ \psi_k: E^k\ra E^i.
$$

By the Telescope Lemma, we have:

\begin{lem}\label{Phi}
    For $0<i<k$, the map $\Phi_{i,k}$ admits an analytic extension to a $d^{k-i}$-to-1 branched covering
from some puzzle piece $F^k\supset E^k$ onto $\hat E^i$.
\end{lem}

\subsection{Travel times} 
  Consider two puzzle pieces $P$ and $Q$ for some GPL or kingdom map $F$.
If $F^l P = Q$, we let $\Time_F (P,Q) =l$
(note that time $l$ is uniquely determined). 
For the ``absolute time'' measured with respect to the initial map $g$, we use notation
 $\Time(P,Q) \equiv \Time_g (P,Q)$. 

   
Let

\ssk \nin $\bullet$ $T_n= \Time (A^n, W^n)$, 
i.e., $ G |A^n = g^{T_n}|A^n$
(the travel time that the king  spends  away from his castle);

\ssk\nin $\bullet$ $t_n= \Time(W^n, V^n)$, i.e., $g_n| W^n = g^{t_n}$;  

\ssk\nin $\bullet$ $s_n = \Time (W^n, W^{n-1})  = t_n+ T_{n-1}$ for $n\geq 1$;  $s_0=t_0=\Time (W^0, V^0)$. 

% \ssk\nin $\bullet$ $N_n$ be the length of the central cascade in the $n$th kingdom.

\begin{lem}\label{t-T-s}
The travel times satisfy the following inequalities:
    $$     t_n\geq T_{n-1}\,; \quad     T_n \geq s_n \, ; \quad    s_n \geq 2 s_{n-1} . $$

\end{lem} 

\begin{pf}
By definition, $g_n(W^n)$ is the first return of $W^n$ to $V^n=A^{n-1}$ under iterates of $G_{n-1}$,
so that $g_n|W^n = G_{n-1}^k|W^n$ for some $k\geq 1$.
Hence 
$$
   g^{t_n}|W^n = g_n|W^n = G_{n-1}^{\circ(k-1)}\circ G_{n-1}|W^n = g^s\circ g^{T_{n-1}}|W^n
$$ 
for some $s\geq 0$,
and the first inequality follows. 

For the second inequality, notice that $T_n$ is the first return time of the critical orbit
to $W^n$ after the first entry to the annulus $V^n\sm W^n$. 
The first entry to $V^n\sm W^n$ occurs at time $\geq t_n$ (since $t_n$ is the first return time of 0 to $V^n$).
Return back to $W^n$ from $V^n\sm W^n$ occurs  at time $\geq T_{n-1}$ (since $T_{n-1}$ is the first moment  $T$
when $f^T(V^n)\cap V^n\not=\emptyset$).

Now the  third inequality follows:
$$
  s_n= t_n+T_{n-1} \geq 2 T_{n-1} \geq  2 s_{n-1}. 
$$

\end{pf} 

\begin{cor}\label{growth}
For any $g$  % (whose dynasty consists of at least $n$ kingdoms), 
we have:   $t_n=\Time(W^n,V^n) \geq 2^{n-1}.$
\end{cor}     


\begin{lem}\label{Time}
    $\Time (W^n,W^{n-2}) \geq \Time (V^n,V^0). $
\end{lem}

\begin{pf}
  We have:   
$$
  \Time (W^n,W^{n-2}) = s_n+s_{n-1}= t_n+T_{n-1}+s_{n-1},
$$
while
$$
   \Time (V^n,V^0) = \Time(A^{n-1}, W^{n-1})+ \Time(W^{n-1},V^0) =
$$ 
$$ 
   = T_{n-1} + s_{n-1}+\dots + s_0 .
$$ 
   Thus, the desired inequality is reduced to:
$$
  t_n \geq s_{n-2}+\dots + s_0.
$$
Now the first two inequalities of Lemma \ref{t-T-s} imply that $t_n\geq s_{n-1}$,
and the last one implies that $s_{n-1} \geq s_{n-2}+\dots + s_0 .$

\end{pf}


  Take some $W^n$, and let $l_0$ be the smallest $l\geq \Time(V^n, V^0)$ such that
$g^l (W^n)\subset W^0$. 

\begin{lem}\label{l}
  $l_0\leq \Time(W^n, W^{n-2})$.
\end{lem}

\begin{pf}
  Let $p=\Time(V^n, V^0)$, $l=\Time(W^n, W^{n-2})$. 
By Lemma \ref{Time}, $l \geq p$. Moreover, $g^l (W^n)= W^{n-2}\subset W^0$.
Hence $l\geq l_0$ by definition of $l_0$. 
\end{pf}

We will now make some combinatorial choices.

Fix some (big) $m$. 
Let  $l_0< l_1< l_2<\dots<l_m $ be the $m$ consecutive return moments of the $\orb W^n$ to $W^0$.
In other words, 
$$
   g^{l_k}(W^n) = h^k (g^{l_0}(W^n)),
$$ 
where $h$ is the generalized renormalization of $g$  on $W^0$.  

Let $n> \log_2 m +5$. 
Then by Corollary \ref{growth}, 
\begin{equation}\label{rel}
       \Time_h (W^{n-2}, V^{n-2}) >   \Time_h (W^{n-3}, V^{n-3}) \geq 
\end{equation}
$$
                \geq \Time_{g_1} (W^{n-3}, V^{n-3}) >  2^{n-5} > m.
$$
or, in the absolute time: 

\begin{equation}\label{absolute}
   l_m - l_0 < \Time(W^{n-3}, V^{n-3}) <  \Time (W^{n-2}, V^{n-2}). 
\end{equation}

Putting this estimate together with  Lemma \ref{l}, we conclude:

\begin{lem}\label{last moment}
  $l_m < \Time (W^n, V^{n-2})$.
\end{lem}

\subsection{Degrees}
  
Let  $O= (g^{l_k} (W^n))_{k=0}^m$.
By Lemma \ref{last moment}, $O$ is contained in  the piece $\TT$ of  $\orb_h W^n$
beginning with $W^n$ and ending with  $V^{n-2}$. 
Let us split $\TT$ into five pieces.
 Namely, let $\TT_i$  be the pieces of $\TT$ between
two consecutive domains, $E^i$ and $E^{i-1}$, of the sequence  
\begin{equation}\label{domains}
                W^n \equiv  E^{2n+1}, E^{2n}, \dots , E^{2n-4} \equiv  V^{n-2}.
\end{equation}
Let $O_i= \TT_i\cap O$.

By (\ref{rel}), each $\TT_i$ has length bigger than  $m$.
Hence at most two of the pieces $O_i$ are non-empty, and so one of them
contains at least $m/2$ elements. Let now $O_i$ stand for such a piece.

%Let $B$ stand for the beginning of the trajectory $\TT_i$ and $E$ stand for the
%end (that is, $B$ and $E$ are appropriate consecutive domains of the string (\ref{domains})).
%Let us define a domain $F \Supset E$ as follows:
%If  $E=W^p$ then $F=V^p$; if $E=V^p$ then $F=\De^p\subset U^{p-1}$.

Let us consider the enlargement $\hat E^{i-1}$ of $E^{i-1}$.
Notice that it is  contained in $W^{n-3}$. 
Let us pull  $\hat E^{i-1}$ back along the $h$-orbit of $W^n$.
It inscribes every domain of this orbit,  $W^n, h(W^n), \dots, h^s(W^n)= E^{i-1}$, 
into  a bigger {\it buffer} domain $F, h(F),\dots, h^s(F) = \hat E^{i-1}$. 

By Lemma \ref{Phi}, we have:

\begin{lem}\label{degree bound}
   The map $h^s: F \ra \hat E^{i-1}$ has degree at most $d^5$.
\end{lem}

Moreover, 

\begin{lem}\label{disjoint}
   The domains $h^k(F)$ enclosing the domains of $O_i$ are pairwise disjoint.
\end{lem}

\begin{pf}
  Otherwise there would be two nested domains $h^k(F) \subset h^s (F)$, $k<s$. 
Let $L = s-k$. Pushing $h^k(F)$ forward to $\hat E^{i-1}$ we see that
$h^L (\hat E^{i-1})\supset \hat E^{i-1}$. All the more, $h^L (W^{n-3}) \supset W^{n-3}$,
so that $L \geq  \Time_h (W^{n-3}, V^{n-3})$.

On the other hand, by  (\ref{absolute}),
 $$
    L < \Time_h (W^{n-3}, V^{n-3}), 
$$
contradiction.
\end{pf}

Let us now consider some domain $\La= g^{l_k}(W^n)\in O_i$, and let $\La' = g^{l_k} (F)$ be its  buffer.
Since there is a biholomorphic push-forward $(\La', \La)\ra ( \hat E^{i-1}, E^{i-1} )$, we have:

\begin{lem}\label{moduli of buffers}
   $ \mod(\La'\sm \La) = \mod (\hat E^{i-1}\sm  E^{i-1}).$
\end{lem} 
  
Let $\Upsilon= g^{-l_k}(V^0)|0$.

\begin{lem}\label{outer degree}
  We have: $W^n\subset \Upsilon \subset V^n$ and 
$$
   \deg(g^{l_k}: \Upsilon\to V^0) \leq d^{2n+m}.
$$
\end{lem}

\begin{pf}
  The first inclusion is trivial. The  second inclusion, $\Upsilon\subset V^n$,
follows from  $l_k\geq l_0\geq \Time(V^n, V^0)$. 

Let us estimate the degree. Let $s= \Time(V^n,V^0)$. 
Then  
$$
     \deg(g^s : V^n \ra V^0)=d^{2n}. 
$$

Let us now consider the first landing map $H$ to $W^0$.
It is easy to see that each component $Q_j$ of the domain of $H$ 
is mapped biholomorphically onto  $W^0$ and, moreover,
 $H:  Q_j\ra W^0$  admits an extension
to a biholomorphic isomorphism $\tl Q_j\ra V^0$. 
Let $\Upsilon_i = g^{l_i}(\Upsilon) $.  Then we have:
$$
  \Upsilon_0  =       H \circ g^s (\Upsilon)
$$
and 
$$
     \Upsilon_{i+1} = H\circ (g | W^0)|\,  \Upsilon_i , \quad i=0,1,\dots, k-1 \leq m-1, 
$$
and the Telescope Lemma concludes the proof. 
\end{pf}

\subsection{Summary}\label{summary} 
We fix an arbitrary $m$ and let $n>\log_2 m +5$. 
Then for any domain $\La =\La_k = g^{l_k}(W^n) \in O_i$, 
the map $\Psi=\Psi_k= g^{l_k}: W^n \ra \La$ admits a holomorphic extension to a branched covering
\begin{equation}\label{3 domains}
    \Psi: (\Upsilon, F, W^n) \ra (V^0, \La', \La)
\end{equation}
such that:

\begin{itemize}

\item [(P1)] $\deg(\Psi: \Upsilon\ra V^0)\leq d^{2n+m}$;

\item[(P2)] $\deg(\Psi: F\ra \La')\leq d^5$;

\end{itemize}

\nin and

\begin{itemize}
\item [(P3)] $\Upsilon\subset V^n$;   

\item[(P4)] $\mod(\La' \sm \La) = \mod (\hat E^{i-1}\sm  E^{i-1})$.
\end{itemize}

\ssk\nin
Moreover, there are at least $m/2$ domains $\La_k$ in the orbit $O_i$, and their buffers $\La_k'$ are 
pairwise disjoint.


\comm{

\begin{lem}\label{once}
All the domains $g^{l_k} (W^n)$, $k=l_0, \dots, m$,
except at most one, are disjoint from  $W^{n-3}$.
\end{lem}

\begin{pf}
  By Lemma \ref{last moment},  all these  domains have depth bigger than the depth of $V^{n-2}$,
all the more, bigger than the depth of $W^{n-3}$.
Hence if some $g^{l_k}(W^n)$ intersects $W^{n-3}$ then it is contained in $W^{n-3}$.
But then by (\ref{absolute}), all the further domains are contained in 
$$ 
  \bigcup_{s=1}^{t_{n-3}-1} g^s (W^{n-3}),
$$
which is disjoint from $W^{n-3}$. 
\end{pf}
end comm}






